\documentclass[11pt,a4paper]{article}
\usepackage[utf8]{inputenc}
\usepackage[T1]{fontenc}
\usepackage{amsmath,amssymb}
\usepackage{booktabs}
\usepackage{multirow}
\usepackage{graphicx}
\usepackage{xcolor}
\usepackage{hyperref}
\usepackage[margin=2.5cm]{geometry}
\usepackage{natbib}
\usepackage{tabularx}
\usepackage{array}

\newcommand{\sys}{OntoCoDec}
\newcommand{\drr}{domain/range violation rate}

\title{\sys{}: Near-Complete Type-Valid Information Extraction\\via Compiled Ontology Constraints}

\author{Anonymous}
\date{}

\begin{document}
\maketitle

%% ============================================================
%% ABSTRACT
%% ============================================================
\begin{abstract}
LLM-based information extraction (IE) systems produce systematic ontology constraint violations---15.4\% of triples on SciERC and 18.7\% on DocRED violate domain/range type constraints.
Schema-informed prompting reduces violations by up to 62\% on simple ontologies but scales poorly to complex ones (only 32\% on DocRED's 96-relation ontology), and can even \emph{increase} violations for specific relation types.
We present \sys{}, a framework that automatically compiles IE ontology schemas into two composable constrained decoding layers: (1)~a \emph{grammar layer} that enforces structural validity and vocabulary compliance via context-free grammar constraints, and (2)~a \emph{semantic layer} that enforces cross-field domain/range constraints through a stateful entity registry with dynamic logit masking.
The two layers compose via bitwise AND of token masks with an enabling dependency: grammar ensures every relation type is valid, enabling semantic lookup of domain/range rules.
On SciERC, \sys{} achieves near-complete violation elimination (domain/range violations 15.4\%$\to$0.93\%, $\downarrow$94\%) at zero F1 cost.
On DocRED, a complex ontology with inherently high constraint tightness, \sys{} achieves predictable partial reduction (19.0\%$\to$10.19\%, $\downarrow$46\%), with 67\% of relation types reaching $<$2\% violation rate.
Residual violations concentrate in a predictable minority of high-tightness relations (Spearman $\rho$=0.52 between constraint tightness and violation rate, $p<$0.001), providing practitioners with a priori estimates of constraint effectiveness.
Grammar-layer enforcement generalizes perfectly across datasets (entity-type and relation-vocabulary violations $\to$0\%) and models (Qwen3-8B, Llama-3.1-8B).
\end{abstract}

%% ============================================================
%% 1. INTRODUCTION
%% ============================================================
\section{Introduction}

Large language models have emerged as powerful zero-shot information extraction systems, capable of extracting structured relational triples from unstructured text without task-specific training.
However, their outputs frequently violate the type constraints specified by the target ontology.
A relation such as \texttt{Feature-of} may require its head entity to be of type \texttt{OtherScientificTerm} and its tail of type \texttt{Method}, yet the model generates type-incompatible entity pairs in over 33\% of instances.

These violations are not merely cosmetic. Downstream knowledge graph population, reasoning engines, and regulatory compliance systems require structurally valid triples.
Post-hoc filtering discards 14--15\% of generated triples, wasting compute without redirecting probability mass toward valid alternatives.
Schema-informed prompting---explicitly describing domain/range constraints in the prompt---provides partial relief but is fundamentally unreliable: on SciERC, it reduces overall violations from 15.4\% to 5.78\% ($\downarrow$62\%), yet the \texttt{Evaluate-for} relation's violation rate \emph{increases} from 24.1\% to 29.9\%, indicating active interference between prompt-level constraint signals and the model's learned type associations.
Moreover, prompt-level intervention scales poorly with ontology complexity: on DocRED (96 relations), reduction drops to only 32\%.

We make three contributions:

\noindent\textbf{(1) Violation Landscape Study.}
We present the first systematic measurement of ontology constraint violations in LLM-based IE across two datasets and two models.
We show that violation rates are driven by \emph{constraint tightness}---the fraction of the entity-type space excluded by a relation's domain/range rules---not ontology size.
SciERC (7 relations) and DocRED (96 relations) exhibit comparable overall violation rates (15.4\% vs.\ 18.7\%) despite a 14$\times$ difference in ontology size.

\noindent\textbf{(2) \sys{} Framework.}
We propose a two-layer constrained decoding framework that automatically compiles IE ontology schemas into composable decoder constraints.
The grammar layer (built on XGrammar) enforces JSON structure, entity-type vocabulary, and relation-type vocabulary via context-free grammar constraints.
The semantic layer (TypeConstraintProcessor) maintains a stateful entity registry during autoregressive generation and dynamically masks logits to enforce domain/range type compatibility.
The layers compose via bitwise AND with an enabling dependency: grammar ensures valid relation vocabulary, enabling semantic domain/range lookup.

\noindent\textbf{(3) Tightness--Violation Correlation.}
We establish that per-relation constraint tightness predicts residual violation rates (Spearman $\rho$=0.52, $p<$0.001 on DocRED), with violations concentrating in a predictable minority of high-tightness relations (top 15 relations account for 76.3\% of all violations).
This provides practitioners with an a priori diagnostic for estimating constraint effectiveness on new ontologies.

%% ============================================================
%% 2. BACKGROUND
%% ============================================================
\section{Background}

\subsection{IE Ontology Constraints}

An IE ontology $\mathcal{O} = (\mathcal{E}, \mathcal{R}, \mathcal{D})$ defines a set of valid entity types $\mathcal{E}$, relation types $\mathcal{R}$, and domain/range constraints $\mathcal{D}: \mathcal{R} \to 2^{\mathcal{E}} \times 2^{\mathcal{E}}$.
For each relation $r \in \mathcal{R}$, the constraint $\mathcal{D}(r) = (\mathcal{E}_h^r, \mathcal{E}_t^r)$ specifies the valid head entity types $\mathcal{E}_h^r \subseteq \mathcal{E}$ and tail entity types $\mathcal{E}_t^r \subseteq \mathcal{E}$.
A triple $(e_h, r, e_t)$ with $\text{type}(e_h) \notin \mathcal{E}_h^r$ or $\text{type}(e_t) \notin \mathcal{E}_t^r$ constitutes a \emph{domain/range violation}.

We define the \emph{constraint tightness} of a relation $r$ as:
\begin{equation}
    \text{tightness}(r) = 1 - \frac{|\mathcal{E}_h^r| \cdot |\mathcal{E}_t^r|}{|\mathcal{E}|^2}
    \label{eq:tightness}
\end{equation}
A tightness of 1.0 means only one entity-type pair is valid; 0.0 means all pairs are valid (no constraint). This metric captures the fraction of the type space excluded by a relation's constraints.

\subsection{Grammar-Constrained Decoding}

Grammar-constrained decoding restricts the output of autoregressive language models to strings accepted by a formal grammar, typically a context-free grammar (CFG).
At each decoding step, a token mask $M_t \in \{0,1\}^{|V|}$ is computed over the vocabulary $V$, zeroing out logits for tokens that would lead to strings outside the grammar's language.
XGrammar compiles grammars into pushdown automata with precomputed token masks, achieving near-zero overhead.

Grammar-constrained decoding can enforce structural constraints (valid JSON, correct field names) and vocabulary constraints (entity types $\in \mathcal{E}$, relation types $\in \mathcal{R}$).
However, it cannot enforce \emph{cross-field} constraints: whether a generated entity type is compatible with the relation type in the same triple requires maintaining state across generation positions, which is beyond CFG expressiveness.

\subsection{Constraint Decidability Taxonomy}

We classify IE ontology constraints by prefix-decidability during autoregressive generation:
\begin{itemize}
    \item \textbf{Type (a) -- Purely local:} Entity type $\in \mathcal{E}$, relation type $\in \mathcal{R}$. Decidable from the current token position alone. Enforceable by grammar.
    \item \textbf{Type (b) -- Prefix-decidable with state:} Domain/range compatibility. Requires maintaining an entity registry across generation positions. Enforceable by stateful logit masking (semantic layer).
    \item \textbf{Type (c) -- Truly non-local:} Forward coreference (entity mentioned later constrains earlier type assignment). Requires future context; out of scope for forward-only decoding.
\end{itemize}

%% ============================================================
%% 3. VIOLATION LANDSCAPE STUDY
%% ============================================================
\section{Violation Landscape Study}
\label{sec:landscape}

Before proposing a solution, we systematically measure the extent and distribution of ontology constraint violations in LLM-based zero-shot IE.

\subsection{Experimental Setup}

We evaluate Qwen3-8B on two IE benchmarks with structured ontology definitions:

\noindent\textbf{SciERC} contains 551 scientific abstracts annotated with 6 entity types and 7 relation types. The ontology is compact but several relations have tight domain/range constraints.

\noindent\textbf{DocRED} contains 998 Wikipedia documents annotated with 6 entity types and 96 relation types drawn from Wikidata. The ontology is large and heterogeneous, with wide variance in constraint tightness (mean tightness 0.842, range 0.50--0.97).

We measure four violation types: (1)~\emph{entity-type violations} (ET): generated entity type $\notin \mathcal{E}$; (2)~\emph{relation-vocabulary violations} (RV): generated relation type $\notin \mathcal{R}$; (3)~\emph{domain/range violations} (DR): entity type incompatible with relation's domain/range constraints; (4)~\emph{cross-field consistency violations} (CF): same entity assigned different types across relations.

\subsection{Unconstrained Baseline Results}

Table~\ref{tab:landscape} summarizes the violation landscape.
Both datasets exhibit substantial DR rates (15.4\% SciERC, 18.7\% DocRED) despite the 14$\times$ difference in ontology size, confirming that violation rates are driven by constraint structure rather than ontology scale.

Schema-informed prompting---embedding explicit domain/range constraint descriptions in the prompt---reduces DR by 62\% on SciERC but only 32\% on DocRED, demonstrating that prompt-level intervention scales poorly with ontology complexity.
On SciERC, three relation types retain $>$10\% DR even with schema-informed prompting (Compare: 27.6\%, Evaluate-for: 29.9\%, Feature-of: 11.3\%).
The \texttt{Evaluate-for} violation rate \emph{increases} with schema-informed prompting (24.1\%$\to$29.9\%), indicating prior-constraint conflict.

\begin{table}[t]
\centering
\small
\caption{Violation landscape: unconstrained LLM-based IE (Qwen3-8B, zero-shot). DR = domain/range violation rate. Schema = schema-informed prompt.}
\label{tab:landscape}
\begin{tabular}{lcccc}
\toprule
\textbf{Dataset} & $|\mathcal{R}|$ & \textbf{DR (\%)} & \textbf{Schema DR (\%)} & \textbf{$\Delta$ DR} \\
\midrule
SciERC  & 7  & 15.4 & 5.78 & $\downarrow$62\% \\
DocRED  & 96 & 18.7 & 12.8 & $\downarrow$32\% \\
\bottomrule
\end{tabular}
\end{table}

\subsection{Per-Relation Violation Distribution}

Violations are not uniformly distributed across relation types.
On SciERC, \texttt{Used-for} has only 0.8\% DR while \texttt{Compare} reaches 42.4\%.
The key differentiator is constraint tightness: \texttt{Compare} requires both arguments to be of type \texttt{Method}, excluding $>$97\% of the entity-type space, while \texttt{Used-for} accepts a broad range of type pairs.

On DocRED (Section~\ref{sec:tightness}), we formally quantify this relationship: Spearman $\rho$=0.52 ($p<$0.001) between constraint tightness and per-relation DR, with a clear monotonic trend across tightness buckets.

%% ============================================================
%% 4. ONTOCODEC FRAMEWORK
%% ============================================================
\section{The \sys{} Framework}
\label{sec:method}

\sys{} automatically compiles an IE ontology schema $\mathcal{O}$ into constrained decoding layers that enforce type validity during autoregressive generation.
The framework consists of three components: an ontology compiler, a grammar constraint layer, and a semantic constraint layer.

\subsection{Ontology Compiler}

The compiler takes as input an ontology definition $\mathcal{O} = (\mathcal{E}, \mathcal{R}, \mathcal{D})$ and produces:
\begin{enumerate}
    \item A context-free grammar $G$ encoding valid JSON structure with entity-type and relation-type vocabularies;
    \item A domain/range lookup table $T: \mathcal{R} \to 2^{\mathcal{E}} \times 2^{\mathcal{E}}$ for runtime type checking.
\end{enumerate}
The compilation is fully automatic---practitioners provide only the ontology definition (entity types, relation types, and domain/range constraints), requiring zero additional code for schema compliance.

\subsection{Grammar Layer}

The grammar layer compiles $G$ into an XGrammar pushdown automaton.
At each decoding step $t$, it computes a mask $M_t^{\text{gram}} \in \{0,1\}^{|V|}$ that allows only tokens leading to valid JSON with:
\begin{itemize}
    \item Correct JSON structure (entity list + relation list with typed fields)
    \item Entity types $\in \mathcal{E}$
    \item Relation types $\in \mathcal{R}$
\end{itemize}
This eliminates entity-type violations (ET$\to$0\%) and relation-vocabulary violations (RV$\to$0\%) by construction, regardless of the model or ontology size.

\subsection{Semantic Layer: TypeConstraintProcessor}

The semantic layer maintains a stateful entity registry $\mathcal{S}$ during generation.
As entities are generated in the entity list portion of the output, $\mathcal{S}$ records each entity's name and type.
When the model generates a relation triple, the processor:
\begin{enumerate}
    \item Identifies the current generation position (head entity type, relation type, or tail entity type);
    \item Looks up the relation type $r$ in the domain/range table $T$;
    \item Computes a mask $M_t^{\text{sem}}$ that allows only entity types compatible with $\mathcal{D}(r)$.
\end{enumerate}

The semantic layer enforces Type~(b) constraints (prefix-decidable with state) that are beyond CFG expressiveness.

\subsection{Layer Composition}

The two layers compose via bitwise AND:
\begin{equation}
    M_t = M_t^{\text{gram}} \wedge M_t^{\text{sem}}
    \label{eq:compose}
\end{equation}

This composition has an \emph{enabling dependency}: the grammar layer ensures every generated relation type $r \in \mathcal{R}$, which is a prerequisite for the semantic layer to look up $\mathcal{D}(r)$.
Without grammar enforcement, the semantic layer encounters unknown relation types for which no domain/range rules exist, rendering it unable to enforce constraints on those triples.

\textbf{Graceful degradation.}
When $M_t^{\text{sem}} \cap M_t^{\text{gram}} = \emptyset$ (the semantic layer would mask all grammar-valid tokens), the system falls back to grammar-only constraints: $M_t = M_t^{\text{gram}}$.
This maintains structural validity as a guaranteed floor.
In practice, this fallback is rare: 1/998 documents on DocRED (0.1\%), 0/551 on SciERC.

\subsection{Decoding Strategies}

We evaluate two strategies:

\noindent\textbf{S1 (Single-pass):} Both grammar and semantic layers are active throughout generation. Recommended as the primary strategy due to simplicity and marginally better metrics.

\noindent\textbf{S2 (Two-pass):} Pass~1 extracts entities with grammar-only constraints; Pass~2 generates relations with an input-dependent grammar compiled from Pass~1 entities.
S2 is an architectural variant demonstrating framework flexibility; it does not improve over S1 on any metric (Section~\ref{sec:exp}).

%% ============================================================
%% 5. EXPERIMENTS
%% ============================================================
\section{Experiments}
\label{sec:exp}

\subsection{Setup}

We evaluate \sys{} on SciERC (551 sentences) and DocRED (998 documents) using Qwen3-8B in a zero-shot setting.
We compare five configurations: \emph{unconstrained} (baseline), \emph{grammar-only} (grammar layer only), \emph{semantic-only} (semantic layer only), \emph{full S1} (both layers, single-pass), and \emph{full S2} (both layers, two-pass).
Cross-model generalization is evaluated with Llama-3.1-8B on SciERC.
F1 is reported as a control variable to verify that constraints do not degrade extraction quality, not as an optimization target.

\subsection{Grammar Layer: Perfect Cross-Dataset Generalization}

Table~\ref{tab:grammar} presents grammar-layer results across datasets.
Grammar-constrained decoding eliminates entity-type violations entirely (ET$\to$0\%) and relation-vocabulary violations to near-zero (RV$\to$0\% on SciERC, 0.03\% on DocRED) across both datasets, scaling from 7 to 96 relation types (14$\times$) with no degradation.
The 5 residual RV violations on DocRED (out of 17,341 triples) likely stem from tokenizer boundary effects on long relation names.

\begin{table}[t]
\centering
\small
\caption{Grammar layer eliminates ET and RV violations across datasets. Grammar-only configuration; violations in constrained output.}
\label{tab:grammar}
\begin{tabular}{lccccc}
\toprule
\textbf{Dataset} & $|\mathcal{R}|$ & \textbf{Triples} & \textbf{ET (\%)} & \textbf{RV (\%)} \\
\midrule
SciERC  & 7  & 1,725  & 0.0 & 0.0 \\
DocRED  & 96 & 17,341 & 0.0 & 0.03 \\
\bottomrule
\end{tabular}
\end{table}

\subsection{SciERC: Near-Complete Violation Elimination}

Table~\ref{tab:scierc} presents the full ablation on SciERC.
Full S1 reduces DR from 15.4\% to 0.93\% ($\downarrow$94\%) while maintaining F1 (8.36\% vs.\ 8.23\% unconstrained; the +0.13pp difference is not statistically significant).

The enabling dependency between layers is quantified by the \emph{enabling ratio}: the DR reduction achieved by full S1 divided by that of semantic-only.
On SciERC, full S1 achieves 14.4pp DR reduction vs.\ semantic-only's 3.7pp, yielding an enabling ratio of 3.88$\times$.
Semantic-only achieves only 24.2\% DR reduction because $\sim$1.4\% of triples use out-of-vocabulary relation types for which no domain/range rules can be looked up.

Full S2 shows no improvement over S1: DR 0.97\% (vs.\ 0.93\%), CF 8.03\% (vs.\ 9.90\%), R-F1 8.03\% (vs.\ 8.36\%).
S2 is positioned as an architectural variant, not a recommended strategy.

\begin{table}[t]
\centering
\small
\caption{SciERC constrained decoding ablation (Qwen3-8B, 551 sentences). DR = domain/range violation rate (\%), CF = cross-field consistency violation rate (\%), R-F1 = relation F1 (\%).}
\label{tab:scierc}
\begin{tabular}{lcccccc}
\toprule
\textbf{Config} & \textbf{ET} & \textbf{RV} & \textbf{DR} & \textbf{CF} & \textbf{R-F1} & \textbf{DR $\Delta$} \\
\midrule
Unconstrained   & $>$0 & 1.4 & 15.4 & $<$0.1 & 8.23 & --- \\
Grammar-only    & 0.0  & 0.0 & $\sim$15.4 & $<$0.1 & --- & $\sim$0\% \\
Semantic-only   & ---  & --- & 11.7 & 2.7  & --- & $\downarrow$24\% \\
Full S1         & 0.0  & 0.0 & 0.93 & 9.9  & 8.36 & $\downarrow$94\% \\
Full S2         & 0.0  & 0.0 & 0.97 & 8.0  & 8.03 & $\downarrow$94\% \\
\bottomrule
\end{tabular}
\end{table}

\subsection{DocRED: Predictable Partial Reduction}

Table~\ref{tab:docred} presents results on DocRED.
Full S1 reduces DR from 19.0\% to 10.19\% ($\downarrow$46.4\%).
While the absolute reduction is lower than SciERC's 94\%, this is predicted by DocRED's higher mean constraint tightness (0.842): tighter constraints force the semantic mask to exclude high-probability tokens, making violations harder to prevent without quality degradation.

The enabling dependency is preserved: semantic-only achieves only 18.4\% DR reduction (enabling ratio 2.52$\times$).
This ratio is lower than SciERC's 3.88$\times$ because DocRED's unconstrained RV is much higher (28.9\% vs.\ 1.4\%), meaning the semantic layer is blind to nearly one-third of triples without grammar enforcement.
This confirms that the enabling dependency becomes \emph{more necessary} as ontology complexity increases.

F1 is preserved: R-F1 8.17\% (full S1) vs.\ 8.20\% (unconstrained), confirming F1-neutral enforcement.

\begin{table}[t]
\centering
\small
\caption{DocRED constrained decoding ablation (Qwen3-8B, 998 documents). ET/RV/DR/CF as violation rates (\%).}
\label{tab:docred}
\begin{tabular}{lcccccc}
\toprule
\textbf{Config} & \textbf{ET} & \textbf{RV} & \textbf{DR} & \textbf{CF} & \textbf{R-F1} & \textbf{DR $\Delta$} \\
\midrule
Unconstrained   & 7.5  & 28.9 & 19.0 & $<$0.1 & 8.20 & --- \\
Grammar-only    & 0.0  & 0.03 & ---  & $<$0.1 & 8.24 & --- \\
Semantic-only   & 7.3  & 28.7 & 15.5 & 2.5  & 8.00 & $\downarrow$18\% \\
Full S1         & 0.0  & 0.03 & 10.19 & 2.6  & 8.17 & $\downarrow$46\% \\
\bottomrule
\end{tabular}
\end{table}

\subsection{Per-Relation DR Distribution on DocRED}
\label{sec:bimodal}

The aggregate 10.19\% DR on DocRED masks a highly skewed per-relation distribution.
Table~\ref{tab:bimodal} summarizes the distribution across 82 relation types with $\geq$5 triples.

\begin{table}[t]
\centering
\small
\caption{DocRED full S1 per-relation DR distribution (82 relations with $\geq$5 triples). Violations are concentrated in a small minority of high-DR relations.}
\label{tab:bimodal}
\begin{tabular}{lcc}
\toprule
\textbf{DR Range} & \textbf{\# Relations} & \textbf{Cumulative \%} \\
\midrule
DR = 0\%      & 18 & 22\% \\
0\% $<$ DR $<$ 2\%   & 37 & 67\% \\
2\%--10\%     & 12 & 82\% \\
10\%--20\%    & 7  & 90\% \\
20\%--40\%    & 6  & 98\% \\
$>$ 40\%      & 8  & 100\% \\
\midrule
\multicolumn{3}{l}{\textit{Summary statistics:}} \\
\multicolumn{2}{l}{Median DR} & 0.8\% \\
\multicolumn{2}{l}{Mean DR} & 10.2\% \\
\multicolumn{2}{l}{Gini coefficient} & 0.82 \\
\multicolumn{2}{l}{Top 15 relations' share of violations} & 76.3\% \\
\bottomrule
\end{tabular}
\end{table}

67\% of relations achieve DR $<$ 2\% (near-complete compliance), while the top 15 high-violation relations contribute 76.3\% of all violations.
The mean/median ratio of 12.75$\times$ and Gini coefficient of 0.82 confirm extreme concentration.
This is not random noise but a systematic consequence of ontology constraint structure, as Section~\ref{sec:tightness} demonstrates.

\subsection{Tightness--Violation Correlation}
\label{sec:tightness}

We compute the Spearman rank correlation between constraint tightness (Eq.~\ref{eq:tightness}) and per-relation DR on DocRED.
Table~\ref{tab:tightness} shows a clear monotonic trend across tightness buckets.

\begin{table}[t]
\centering
\small
\caption{DocRED per-relation DR by constraint tightness bucket (full S1). Higher tightness $\to$ higher residual DR. Spearman $\rho$=0.52, $p<$0.001.}
\label{tab:tightness}
\begin{tabular}{lccc}
\toprule
\textbf{Tightness Bucket} & \textbf{\# Rels} & \textbf{Agg.\ DR (\%)} \\
\midrule
$<$ 0.70          & 2  & 0.7 \\
0.70--0.79        & 18 & 9.0 \\
0.80--0.84        & 23 & 7.0 \\
0.85--0.89        & 28 & 12.9 \\
0.90--0.94        & 15 & 19.3 \\
0.95--1.00        & 4  & 30.6 \\
\bottomrule
\end{tabular}
\end{table}

The Spearman $\rho$=0.52 ($n$=82, $p<$0.001) confirms a statistically significant positive correlation.
While the correlation is moderate in absolute terms, the bucket-level monotonic trend and extreme violation concentration (76.3\% in top 15 relations) provide strong practical utility: practitioners can estimate constraint effectiveness from the ontology definition alone, before any model inference.

The mechanism is intuitive: high-tightness relations exclude most of the entity-type space, forcing the semantic mask to block high-probability tokens.
When the model's distributional prior strongly favors invalid type assignments (e.g., \texttt{P6\_head\_of\_government} with tightness 0.889 and DR 84.6\%), even per-token logit masking cannot fully redirect probability mass without degrading generation coherence.

Illustrative extreme cases include \texttt{P6\_head\_of\_government} (tightness 0.889, DR 84.6\%, $n$=13) and \texttt{P582\_end\_time} (tightness 0.917, DR 43.5\%, $n$=186), both of which have very narrow valid type pairs.
In contrast, \texttt{P17\_country} and similar low-tightness relations achieve near-zero DR.

\subsection{Cross-Model Generalization}
\label{sec:crossmodel}

To verify that our findings are not model-specific, we replicate the SciERC experiments with Llama-3.1-8B.
Table~\ref{tab:crossmodel} compares violation rates across models.

\begin{table}[t]
\centering
\small
\caption{Cross-model comparison on SciERC. Grammar enforcement (ET/RV$\to$0\%) and DR reduction generalize across models. $\rho$ = per-relation DR rank Spearman correlation between models.}
\label{tab:crossmodel}
\begin{tabular}{llcccc}
\toprule
\textbf{Model} & \textbf{Config} & \textbf{ET} & \textbf{RV} & \textbf{DR} & \textbf{R-F1} \\
\midrule
\multirow{3}{*}{Qwen3-8B}
& Unconstrained & $>$0 & 1.4 & 15.4 & 8.23 \\
& Grammar       & 0.0  & 0.0 & ---  & --- \\
& Full S1       & 0.0  & 0.0 & 0.93 & 8.36 \\
\midrule
\multirow{3}{*}{Llama-3.1-8B}
& Unconstrained & 0.6  & 2.5 & 14.5 & 6.50 \\
& Grammar       & 0.0  & 0.0 & 14.0 & --- \\
& Full S1       & 0.0  & 0.0 & 1.2  & 6.80 \\
\bottomrule
\multicolumn{6}{l}{\textit{Cross-model per-relation DR Spearman} $\rho$=0.89} \\
\end{tabular}
\end{table}

Grammar enforcement achieves ET/RV$\to$0\% on both models, confirming that grammar constraints are model-agnostic (they depend only on the tokenizer vocabulary mapping, automatically adapted by XGrammar).

Full S1 reduces DR from 14.5\% to 1.2\% on Llama ($\downarrow$92\%), comparable to Qwen's $\downarrow$94\%.
The cross-model per-relation DR Spearman $\rho$=0.89 indicates that the \emph{per-relation violation susceptibility is a property of the ontology constraint structure, not the generating model}---both models violate the same relations at similar rates, driven by constraint tightness.

\subsection{Cross-Field Consistency Analysis}
\label{sec:cf}

Constrained decoding introduces a trade-off: while DR violations decrease, cross-field consistency (CF) violations increase.
CF violations occur when the same entity is assigned different types across different relations---each individual triple remains type-valid, but cross-triple consistency is broken.

Table~\ref{tab:cf} presents the CF analysis.

\begin{table}[t]
\centering
\small
\caption{Cross-field consistency (CF) analysis. CF increases with semantic enforcement but at a favorable trade-off ratio (DR reduction per CF increase). EVR = entity-vs-relation type drift (dominant CF sub-type).}
\label{tab:cf}
\begin{tabular}{lccccc}
\toprule
\textbf{Dataset} & \textbf{Config} & \textbf{CF (\%)} & \textbf{DR (\%)} & \textbf{DR:CF} \\
\midrule
\multirow{3}{*}{SciERC}
& Unconstrained & $<$0.1 & 15.4 & --- \\
& Grammar-only  & 0.08   & $\sim$15.4 & --- \\
& Full S1       & 9.9    & 0.93 & 1.5:1 \\
\midrule
\multirow{3}{*}{DocRED}
& Unconstrained & $<$0.1 & 19.0 & --- \\
& Grammar-only  & 0.07   & ---  & --- \\
& Full S1       & 2.6    & 10.19 & 3.6:1 \\
\bottomrule
\end{tabular}
\end{table}

\textbf{CF is 100\% attributable to the semantic layer.}
Grammar-only CF rates ($\leq$0.08\%) are identical to unconstrained baselines, confirming that grammar enforcement introduces no CF.
All CF increase comes from the semantic layer's per-relation independent type enforcement: the semantic mask selects the locally optimal type for each relation's domain/range constraints, without maintaining global entity-type consensus.

\textbf{Root cause: type system ambiguity.}
On SciERC, 90\% of CF violations are entity-vs-relation type drift (EVR), where the same entity receives different types in the entity list vs.\ relation triples.
The dominant conflict patterns---Generic$\leftrightarrow$OtherScientificTerm (60 cases), Material$\leftrightarrow$OtherScientificTerm (28), Material$\leftrightarrow$Method (20)---reflect genuine semantic boundary ambiguity in SciERC's type system.

The cross-dataset comparison provides strong evidence for this hypothesis: DocRED CF (2.6\%) is 3.8$\times$ lower than SciERC CF (9.9\%).
DocRED's entity types (PER/ORG/LOC/TIME/NUM/MISC) have far clearer semantic boundaries than SciERC's (Generic/OtherScientificTerm/Material/Method), confirming that \emph{CF rate is driven by type system ambiguity}.
On DocRED, the dominant CF pattern is ORG$\leftrightarrow$LOC (23.1\% of CF), corresponding to the well-known NER ambiguity of geopolitical entities.

\textbf{Severity distinction.}
DR violations produce structurally invalid triples unusable by downstream systems (hard constraint).
CF violations produce individually valid triples with cross-triple inconsistency (soft preference).
The trade-off is favorable: on DocRED, each 1pp CF increase buys 3.57pp DR reduction.

\subsection{Inference Latency}

Table~\ref{tab:latency} reports inference overhead.
Full S1 incurs 1.37--1.39$\times$ overhead across datasets---the grammar layer adds $\sim$13\% (XGrammar's optimized C++ pushdown automaton) and the semantic layer adds $\sim$20\% (Python-based per-token logit masking).
The overhead is per-token and independent of document length.
Full S2 incurs 2.17$\times$ overhead due to two forward passes, with no quality benefit.

\begin{table}[t]
\centering
\small
\caption{Inference latency overhead (tokens/second relative to unconstrained).}
\label{tab:latency}
\begin{tabular}{lcccc}
\toprule
\textbf{Dataset} & \textbf{Grammar} & \textbf{Semantic} & \textbf{Full S1} & \textbf{Full S2} \\
\midrule
SciERC & 1.13$\times$ & 1.33$\times$ & 1.37$\times$ & 2.17$\times$ \\
DocRED & 1.14$\times$ & --- & 1.39$\times$ & --- \\
\bottomrule
\end{tabular}
\end{table}

%% ============================================================
%% 6. RELATED WORK
%% ============================================================
\section{Related Work}

\textbf{Constrained decoding for structured generation.}
Grammar-constrained decoding has been applied to JSON generation, code synthesis, and SQL queries.
XGrammar and related systems compile formal grammars into efficient token masks.
Recent work on code generation applies type-checking during decoding in programming language contexts.
\sys{} extends this paradigm to IE ontology constraints, introducing semantic-layer enforcement for cross-field constraints beyond CFG expressiveness.

\textbf{LLM-based information extraction.}
Zero-shot and few-shot LLM-based IE has shown strong results on standard benchmarks.
However, prior work focuses primarily on F1 metrics and largely ignores ontology compliance.
To our knowledge, no prior work systematically measures or addresses ontology constraint violations in LLM-generated extractions.

\textbf{Ontology-aware IE.}
Traditional IE systems incorporate ontology constraints through training-time loss functions, post-hoc filtering, or ILP-based inference.
Post-hoc filtering cannot redirect probability mass toward valid outputs and wastes compute on invalid generations.
Training-time approaches require task-specific supervision.
\sys{} operates at decoding time with zero training, automatically compiled from ontology definitions.

\textbf{Constraint satisfaction in NLP.}
Lexically constrained decoding enforces keyword inclusion during generation.
Neurologic and A*-based methods handle more complex constraints.
These approaches target surface-level lexical constraints rather than semantic type constraints.
\sys{}'s semantic layer addresses a different constraint class---cross-field type compatibility---requiring stateful tracking during generation.

%% ============================================================
%% 7. CONCLUSION
%% ============================================================
\section{Conclusion}

We presented \sys{}, a framework for enforcing ontology type constraints during LLM-based information extraction via compiled constrained decoding.
Our violation landscape study reveals that LLM-based IE systems produce systematic, non-uniformly distributed ontology violations driven by constraint tightness.
\sys{}'s two-layer architecture---grammar for structural/vocabulary compliance, semantic for cross-field domain/range enforcement---composes via bitwise AND with an enabling dependency, achieving near-complete violation elimination on tightly-constrained ontologies (94\% DR reduction on SciERC) and predictable partial reduction on complex ontologies (46\% on DocRED), all at zero F1 cost and $<$1.4$\times$ latency overhead.
Grammar-layer enforcement generalizes perfectly across datasets and models.

The tightness--violation correlation (Spearman $\rho$=0.52) provides practitioners with an a priori diagnostic for estimating constraint effectiveness from ontology definitions alone, without model inference.

\textbf{Limitations.}
Cross-field consistency violations are an inherent trade-off of per-relation independent type enforcement (9.9\% on SciERC, 2.6\% on DocRED), with CF rate driven by type system ambiguity.
Future work should explore global entity-type consensus mechanisms---e.g., a pre-decoding entity-type resolution pass---to maintain cross-relation consistency.
Our evaluation covers two datasets and two models; broader validation across ontologies, languages, and model scales is needed.
The current implementation uses a Python-based semantic layer; a C++ implementation would reduce the 20\% semantic overhead.

\bibliographystyle{plainnat}

\end{document}
